\section*{Course Purpose and Outcomes}
CIVL 2050 is the introductory course in fluid and flow phenomena applicable to civil and environmental engineering problems. It is also the laboratory course in fluids and flow phenomena. The course emphasizes real-fluid behavior, engineering analysis, experimental measurement, and interpretation of observations in the context of theory.

\textbf{Students who successfully complete this course shall be able to:}
\begin{enumerate}
\item Analyze pressures exerted by fluids at rest and in uniform flow.
\item Apply conservation of mass, energy, and momentum to a control volume.
\item Analyze energy losses in pipe flows.
\item Use dimensionless numbers to evaluate pipe friction, pump performance, and model/prototype relationships.
\item Acquire skills in analyzing and solving engineering problems.
\item Make measurements of flow phenomena.
\item Apply hydrostatics, conservation of mass, energy, and momentum to control-volume problems.
\item Write succinct reports on experiments and discuss observations in the context of theory.
\end{enumerate}

\subsection*{Course Assessment Measures}
\begin{tabularx}{\textwidth}{@{}L{0.20\textwidth} L{0.25\textwidth} X@{}}
\toprule
Assessment & Due date / timing & Learning outcome numbers \\
\midrule
Exams & As specified below & 1, 2, 3, 4, 7 \\
Laboratories & As specified below & 1, 2, 3, 4, 5, 6, 7, 8 \\
Homework & As specified below & 1, 2, 3, 4, 7 \\
\bottomrule
\end{tabularx}
